\textbf{AUTORIZACIÓN PARA LA DIGITALIZACIÓN, DEPÓSITO Y DIVULGACIÓN EN RED DE PROYECTOS FIN DE GRADO, FIN DE MÁSTER, TESINAS O MEMORIAS DE BACHILLERATO}
\vspace*{1em}

\textbf{\textit{1º. Declaración de la autoría y acreditación de la misma.}}
\vspace*{1em}
 
El autor D.\underline{\hspace{6cm}} \textbf{DECLARA} ser el titular de los derechos de 
propiedad intelectual de la obra: \underline{\hspace{6cm}}, que ésta es una obra original, 
y que ostenta la condición de autor en el sentido que otorga la Ley de Propiedad Intelectual.
 
\vspace*{1em}
\textbf{\textit{2º. Objeto y fines de la cesión.}}
\vspace*{1em}

Con el fin de dar la máxima difusión a la obra citada a través del Repositorio 
institucional de la Universidad, el autor \textbf{CEDE} a la Universidad Pontificia Comillas,
 los derechos de digitalización, de archivo, de reproducción, de distribución y 
 de forma gratuita y no exclusiva, por el máximo plazo legal y con ámbito universal,
de comunicación pública, incluido el derecho de puesta a disposición electrónica, 
tal y como se describen en la Ley de Propiedad Intelectual. El derecho de 
transformación se cede a los únicos efectos de lo dispuesto en la letra a) 
del apartado siguiente.

\vspace*{1em}
\textbf{\textit{3º. Condiciones de la cesión y acceso}} 
\vspace*{1em}

Sin perjuicio de la titularidad de la obra, que sigue correspondiendo a su autor, 
la cesión de derechos contemplada en esta licencia habilita para:
 
\begin{enumerate}[(a)]
	\item Transformarla con el fin de adaptarla a cualquier tecnología que permita incorporarla a internet y hacerla accesible; incorporar metadatos para realizar el registro de la obra e incorporar $``$marcas de agua$"$  o cualquier otro sistema de seguridad o de protección. 
	\item Reproducirla en un soporte digital para su incorporación a una base de datos electrónica, incluyendo el derecho de reproducir y almacenar la obra en servidores, a los efectos de garantizar su seguridad, conservación y preservar el formato.
	\item  Comunicarla, por defecto, a través de un archivo institucional abierto, accesible de modo libre y gratuito a través de internet. 
	\item Cualquier otra forma de acceso (restringido, embargado, cerrado) deberá solicitarse expresamente y obedecer a causas justificadas.
	\item  Asignar por defecto a estos trabajos una licencia Creative Commons. 
	\item Asignar por defecto a estos trabajos un HANDLE (URL \textit{persistente}).
\end{enumerate}

\vspace*{1em}
\textbf{\textit{4º. Derechos del autor.}}
\vspace*{1em}

El autor, en tanto que titular de una obra tiene derecho a: 

\begin{enumerate}[(a)]
	\item  Que la Universidad identifique claramente su nombre como autor de la misma 
	\item  Comunicar y dar publicidad a la obra en la versión que ceda y en otras posteriores a través de cualquier medio. 
	\item  Solicitar la retirada de la obra del repositorio por causa justificada.  
	\item  Recibir notificación fehaciente de cualquier reclamación que puedan formular terceras personas en relación con la obra y, en particular, de reclamaciones relativas a los derechos de propiedad intelectual sobre ella. 
\end{enumerate}

\vspace*{1em}
\textbf{\textit{5º. Deberes del autor.}} 
\vspace*{1em}

\begin{enumerate}[(a)]
	\item El autor se compromete a: 
	\item Garantizar que el compromiso que adquiere mediante el presente escrito no infringe ningún derecho de terceros, ya sean de propiedad industrial, intelectual o cualquier otro. 
	\item Garantizar que el contenido de las obras no atenta contra los derechos al honor, a la intimidad y a la imagen de terceros.
	\item Asumir toda reclamación o responsabilidad, incluyendo las indemnizaciones por daños, que pudieran ejercitarse contra la Universidad por terceros que vieran infringidos sus derechos e intereses a causa de la cesión. 
	\item Asumir la responsabilidad en el caso de que las instituciones fueran condenadas por infracción de derechos derivada de las obras objeto de la cesión.
\end{enumerate}

\vspace*{1em}
\textbf{\textit{6º. Fines y funcionamiento del Repositorio Institucional.}}
\vspace*{1em}

La obra se pondrá a disposición de los usuarios para que hagan de ella un uso justo y respetuoso con los derechos del autor, según lo permitido por la legislación aplicable, y con fines de estudio, investigación, o cualquier otro fin lícito. Con dicha finalidad, la Universidad asume los siguientes deberes y se reserva las siguientes facultades: 

\begin{itemize}
	\item La Universidad informará a los usuarios del archivo sobre los usos permitidos, y no garantiza ni asume responsabilidad alguna por otras formas en que los usuarios hagan un uso posterior de las obras no conforme con la legislación vigente. El uso posterior, más allá de la copia privada, requerirá que se cite la fuente y se reconozca la autoría, que no se obtenga beneficio comercial, y que no se realicen obras derivadas.
	\item La Universidad no revisará el contenido de las obras, que en todo caso permanecerá bajo la responsabilidad exclusive del autor y no estará obligada a ejercitar acciones legales en nombre del autor en el supuesto de infracciones a derechos de propiedad intelectual derivados del depósito y archivo de las obras. El autor renuncia a cualquier reclamación frente a la Universidad por las formas no ajustadas a la legislación vigente en que los usuarios hagan uso de las obras.
	\item La Universidad adoptará las medidas necesarias para la preservación de la obra en un futuro. 
	\item La Universidad se reserva la facultad de retirar la obra, previa notificación al autor, en supuestos suficientemente justificados, o en caso de reclamaciones de terceros. 
\end{itemize}

\vspace*{2em}

Madrid, a \ldots \ldots  de \ldots \ldots \ldots \ldots \ldots \ldots de \ldots \ldots 

\vspace*{2em}
{\large ACEPTA}
\vspace*{2em}

Fdo.: \ldots \ldots \ldots \ldots \ldots \ldots
\vspace*{2em}

Motivos para solicitar el acceso restringido, cerrado o embargado del trabajo en el Repositorio Institucional:

\begin{tcolorbox}[height fill, sharp corners, colback=white]
\end{tcolorbox}
